\documentclass[12pt]{article}
\usepackage{amsmath,amssymb}

\begin{document}
\section*{{2020 AIME I Problems/Problem 13}}
Source: AoPS Wiki

\subsection*{Solution}
Points are defined as shown (Actually Point G is a point on line BC such that $\triangle AGD$ has circumcenter F, while Point H is a point on line BC such that $\triangle AHD$ has circumcenter E). Since $\angle{AFE} = \angle{AGD}$ , and $\angle{AEF} = \angle{AHD}$ , it is a breakthrough to get $\triangle AFE \sim \triangle AGH$ . We'd like to compare altitudes to help us compare their areas. Now, note that $AD/2$ is the altitude of $\triangle AFE$ , and Stewart's theorem implies that $AD/2 = \frac{\sqrt{18}}{2}$ . Similarly, the altitude of $\triangle AGH$ is the altitude of $\triangle ABC$ , or $\frac{3\sqrt{7}}{2}$ . However, by the symmetry axis $BE$ on quadrilateral $ABHE$ , we have $BH = AB = 4$ , and similarly $CG = CA = 6$ , so $GB = HC = 1$ , and therefore $[AGH] = [ABC]$ . From here, we get that the area of $\triangle AFE$ is $\frac{15\sqrt{7}}{4}*(\frac{\frac{3\sqrt{7}}{2}}{\frac{\sqrt{18}}{2}})^2=\frac{15\sqrt{7}}{14} \implies \boxed{036}$ , by similarity. ~awang11 ~gougutheorem

\end{document}
